\section{Ánh xạ yêu cầu nghiệp vụ sang yêu cầu xử lý dữ liệu}
\label{sec:3-mapping}

Mục này nối các yêu cầu nghiệp vụ ở Mục~\ref{sec:3-br} với các thành phần xử lý dữ liệu sẽ được thiết kế ở Chương~\ref{chap:4-design}. Bảng~\ref{tab:3-traceability} là ma trận truy vết: một dấu \checkmark{} nghĩa là thành phần đó tham gia trực tiếp vào việc đáp ứng yêu cầu.

{\small
\begin{longtblr}[
  caption = {Ma trận truy vết giữa yêu cầu nghiệp vụ và thành phần xử lý dữ liệu},
  label = {tab:3-traceability},
]{
  width = \linewidth,
  rowhead = 1,
  colspec = {X[3,l]X[1,c]X[1,c]X[1,c]X[1.3,c]X[1,c]X[1,c]X[1.3,c]}, rows = {abovesep=2pt, belowsep=2pt},
  hlines, vlines, row{1} = {font=\bfseries\footnotesize, bg=gray!20},
}
Yêu cầu & Kafka & Spark stream & Cassandra & Iceberg (bronze/gold) & dbt & Trino & FastAPI / web \\
BR1 Tập trung hóa giám sát & \checkmark & \checkmark & \checkmark & \checkmark & & \checkmark & \checkmark \\
BR2 Nhận biết điều kiện đồng ruộng & \checkmark & \checkmark & \checkmark & & & & \checkmark \\
BR3 Liên tục giám sát & \checkmark & \checkmark & \checkmark & \checkmark & \checkmark & & \\
BR4 Dữ liệu đáng tin cậy & & \checkmark & \checkmark & \checkmark & \checkmark & \checkmark & \checkmark \\
BR5 Quản lý dựa trên dữ liệu & & & & \checkmark & \checkmark & \checkmark & \checkmark \\
BR6 Đánh giá lịch sử & & & & \checkmark & \checkmark & \checkmark & \checkmark \\
BR7 Thông tin dễ tiếp cận & & & \checkmark & \checkmark & & \checkmark & \checkmark \\
BR8 Ứng phó sớm với rủi ro & \checkmark & \checkmark & \checkmark & \checkmark & & & \checkmark \\
\end{tblr}
\end{table}

\subsection{Cách đọc ma trận}
\label{sec:3-mapping-read}

Ma trận cho thấy ba nhóm yêu cầu, tương ứng với ba đường dữ liệu của hệ thống:

\begin{itemize}
  \item \textbf{Đường gần thời gian thực} (Kafka, Spark streaming, Cassandra): BR2 và BR3 dựa chủ yếu vào đường này, vì chúng cần dữ liệu mới trong khoảng một phút và cần chịu được gián đoạn.
  \item \textbf{Đường phân tích} (Iceberg, dbt, Trino): BR5 và BR6 chỉ dùng đường này, vì chúng làm việc trên dữ liệu tổng hợp nhiều ngày đến nhiều tháng.
  \item \textbf{Yêu cầu xuyên suốt}: BR1, BR4 và BR8 đi qua cả hai đường. Với BR8, luật cảnh báo chạy trong Spark, cảnh báo và trạng thái của chúng nằm trong Cassandra, còn Iceberg nhận dữ liệu thô và các khoảng mất kết nối (\code{ops.device\_gaps}). Lịch sử cảnh báo chưa được đưa vào tầng gold (FR8.6, Bảng~\ref{tab:3-mapping-br4-br8}).
\end{itemize}

\subsection{Yêu cầu xử lý cho BR4 và BR8}
\label{sec:3-mapping-br4-br8}

Bảng~\ref{tab:3-mapping-br4-br8} chi tiết hóa hai yêu cầu xuyên suốt BR4 và BR8 thành các yêu cầu chức năng (FR) và nơi chúng được hiện thực.

{\small
\begin{longtblr}[
  caption = {Ánh xạ yêu cầu chức năng của BR4 và BR8 sang bước xử lý},
  label = {tab:3-mapping-br4-br8},
]{
  width = \linewidth,
  rowhead = 1,
  colspec = {X[0.8,l]X[2.6,l]X[1.6,l]X[2.4,l]}, rows = {abovesep=2pt, belowsep=2pt},
  hlines, vlines, row{1} = {font=\bfseries, bg=gray!20},
}
FR & Yêu cầu & Thành phần & Cách hiện thực \\
FR4.1 & Schema tường minh, cách ly bản ghi lỗi & Spark streaming & \code{StructType} tường minh; bảng \code{ops.quarantine} \\
FR4.2 & Kiểm tra khoảng giá trị, phát hiện bước nhảy & Spark streaming & Ngưỡng \code{min}, \code{max}, \code{max\_jump} trong dữ liệu danh mục \\
FR4.3 & Khử trùng lặp theo \code{(device\_id, metric, measured\_at)} & Spark, dbt & \code{row\_number()} và \code{MERGE} theo khóa \\
FR4.4 & Lưu vết nguồn gốc trên mọi bản ghi & Spark, Iceberg & Cột topic, partition, offset, \code{ingested\_at}, \code{batch\_id} \\
FR4.5 & Tóm tắt chất lượng hằng ngày & dbt, Trino & Model gold \code{quality\_daily} \\
FR4.6 & Bảng điều khiển loại bỏ hoặc đánh dấu dữ liệu lỗi & FastAPI, web & Cột \code{quality\_flag} trả về cùng giá trị \\
FR4.7 & Test dbt chặn làm mới gold khi thất bại & dbt & \code{dbt build} bỏ qua model phía sau test lỗi \\
FR8.1 & Phát hiện vượt ngưỡng theo cây trồng & Spark streaming & Ngưỡng \code{alert\_below}, \code{alert\_above} theo cây trồng \\
FR8.2 & Phát hiện mất kết nối, pin yếu, sóng yếu & Spark streaming & Thời gian im lặng so với chu kỳ báo cáo \\
FR8.3 & Lưu cảnh báo với loại, mức độ, thửa ruộng, trạng thái & Cassandra & Bảng \code{alerts\_by\_field}, \code{alerts\_by\_status}, \code{alert\_by\_id} \\
FR8.4 & Vòng đời open, acknowledged, resolved & Spark, FastAPI & Máy trạng thái thuần trong \code{alerts.py}; \code{PATCH /alerts/\{id\}} \\
FR8.5 & Chống bão cảnh báo (debounce, cooldown) & Spark, Cassandra & Bảng \code{alert\_state\_by\_field} \\
FR8.6 & Danh sách cảnh báo đang mở và lịch sử & Cassandra, FastAPI, web & \code{GET /alerts?status=} và \code{GET /fields/\{id\}/alerts} (lịch sử giữ 90 ngày trong Cassandra); bảng lịch sử cảnh báo ở tầng gold chưa được hiện thực \\
\end{tblr}
\end{table}
